\chaves{Fertilização, Vitro, IVF, Evolucionário, Evolutivo, SPEA2, Multi-Objetivo}

\begin{resumo} 
Problemas de Otimização Multiobjetivo (MOPs) frequentemente apresentam objetivos conflitantes e espaços de busca complexos, impondo desafios aos métodos de otimização tradicionais. O Algoritmo Evolucionário de Força Pareto 2 (SPEA2) é uma abordagem bem estabelecida para lidar com tais problemas, mas pode enfrentar limitações no balanceamento entre convergência e diversidade em cenários complexos. Este artigo propõe o IVF/SPEA2, um algoritmo híbrido que integra o método de Fertilização In Vitro (IVF) ao SPEA2 para aprimorar o desempenho por meio de manipulação genética direcionada. O IVF/SPEA2 foi avaliado utilizando quatro conjuntos de benchmarks amplamente reconhecidos ZDT, DTLZ, WFG e MaF cobrindo uma considerável gama de geometrias de fronteira de Pareto e complexidades de problemas. O desempenho foi medido utilizando a métrica de Distância Geracional Invertida (IGD), com a significância estatística avaliada por meio do teste de soma de postos de Wilcoxon em 100 execuções independentes. Os resultados demonstram que o IVF/SPEA2 alcança melhorias estatisticamente significativas em relação ao SPEA2 canônico em inúmeros problemas bi-objetivo e tri-objetivo, particularmente em cenários envolvendo multimodalidade, engano (deception) e não separabilidade. Ademais, uma análise comparativa abrangente revela que o IVF/SPEA2 consistentemente supera algoritmos populares de estado da arte, incluindo NSGA-II, NSGA-III e MOEA/D, bem como hibridizações mais recentes do SPEA2, como MFOSPEA2 e SPEA2+SDE. O algoritmo proposto demonstra características superiores de convergência e diversidade de soluções em diversos cenários de problemas, estabelecendo sua eficácia tanto contra abordagens clássicas quanto contemporâneas de otimização multiobjetivo. Esses achados validam o potencial do IVF/SPEA2 como um algoritmo memético robusto e adaptável para tarefas complexas de otimização multiobjetivo.
\end{resumo}

